\begin{figure}[H]
\centering
\begin{tikzpicture}[
    scale=1.3,
    dot/.style={circle, inner sep=1.5pt},
    faintdot/.style={dot, fill=black!25},
    sel/.style={dot, fill=black},
    faintedge/.style={line width=0.55pt, draw=black!25},
    seledge/.style={line width=0.85pt, draw=black},
    lab/.style={draw=none, fill=none, inner sep=0pt, font=\small},
]

\begin{scope}[xshift=0cm]
    % vertices
    \foreach \x in {1,2,3}{
    \foreach \y in {1,2,3}{
        \node[faintdot] (L-\x-\y) at (\x,\y) {};
    }
    }

    % select vertices (2,1), (1,2), (2,3), (3,2)
    \node[sel] (Lsel-2-1) at (2,1) {};
    \node[sel] (Lsel-1-2) at (1,2) {};
    \node[sel] (Lsel-2-3) at (2,3) {};
    \node[sel] (Lsel-3-2) at (3,2) {};

    % (faint) edges of CGrid_{3x3} copied from your reference
    \draw[faintedge] (L-1-1) -- (L-1-2);
    \draw[faintedge] (L-1-1) -- (L-2-1);
    \draw[faintedge] (L-1-1) -- (L-2-2);
    \draw[faintedge] (L-1-1) -- (L-2-3);
    \draw[faintedge] (L-1-1) -- (L-3-2);

    \draw[faintedge] (L-1-2) -- (L-1-3);
    \draw[faintedge] (L-1-2) -- (L-2-2);
    \draw[faintedge] (L-1-2) -- (L-2-3);

    \draw[faintedge] (L-1-3) -- (L-2-3);

    \draw[faintedge] (L-2-1) -- (L-2-2);
    \draw[faintedge] (L-2-1) -- (L-3-1);
    \draw[faintedge] (L-2-1) -- (L-3-2);

    \draw[faintedge] (L-2-2) -- (L-2-3);
    \draw[faintedge] (L-2-2) -- (L-3-2);

    \draw[faintedge] (L-2-3) -- (L-3-3);
    \draw[faintedge] (L-3-1) -- (L-3-2);
    \draw[faintedge] (L-3-2) -- (L-3-3);

    \draw[faintedge] (L-1-1) to[bend left=20] (L-1-3);
    \draw[faintedge] (L-2-1) to[bend left=20] (L-2-3);
    \draw[faintedge] (L-3-1) to[bend left=20] (L-3-3);

    \draw[faintedge] (L-1-1) to[bend right=20] (L-3-1);
    \draw[faintedge] (L-1-2) to[bend right=20] (L-3-2);

    \draw[faintedge] (L-1-1) to[bend right=20] (L-3-3);
    \draw[faintedge] (L-1-2) -- (L-3-3);
    \draw[faintedge] (L-1-3) to[bend right=20] (L-3-3);

    \draw[faintedge] (L-2-1) -- (L-3-3);
    \draw[faintedge] (L-2-2) -- (L-3-3);

    % chosen subgraph edges among the 4 chosen vertices
    \draw[seledge] (L-2-1) -- (L-3-2);
    \draw[seledge] (L-2-1) -- (L-2-3);
    \draw[seledge] (L-1-2) -- (L-3-2); 
    \draw[seledge] (L-1-2) -- (L-2-3);

    % corner labels (optional)
    \node[lab, anchor=east] at ($(L-1-1)+(-0.25,0)$) {$(1,1)$};
    \node[lab, anchor=east] at ($(L-1-3)+(-0.25,0)$) {$(1,3)$};
    \node[lab, anchor=west] at ($(L-3-1)+(0.25,0)$) {$(3,1)$};
    \node[lab, anchor=west] at ($(L-3-3)+(0.25,0)$) {$(3,3)$};
\end{scope}

% ---------------------------
% Arrow to sheared grid
% ---------------------------
\draw[->,thick] (3.5,2) to (6.5,2);

% ---------------------------
% Helper: middle grid (shear/shift y by 0.2*(x-1))
% (x,y) -> (x, y + 0.2*(x-1))
% ---------------------------
\begin{scope}[xshift=6cm]
    % vertices
    \foreach \x in {1,2,3}{
    \foreach \y in {1,2,3}{
        \pgfmathsetmacro{\dx}{0.2*(\y-1)} % shift x depending on y
        \pgfmathsetmacro{\dy}{0.2*(\x-1)} % shift y depending on x
        \node[faintdot] (M-\x-\y) at (\x+\dx,\y+\dy) {};
    }
    }

    % select vertices
    \node[sel] (Msel-2-1) at (2,1+0.2) {};
    \node[sel] (Msel-1-2) at (1+0.2,2) {};
    \node[sel] (Msel-2-3) at (2+0.4,3+0.2) {};
    \node[sel] (Msel-3-2) at (3+0.2,2+0.4) {};

    % faint edges: same adjacency, just using shifted coordinates
    \draw[faintedge] (M-1-1) -- (M-1-2);
    \draw[faintedge] (M-1-1) -- (M-2-1);
    \draw[faintedge] (M-1-1) -- (M-2-2);
    \draw[faintedge] (M-1-1) -- (M-2-3);
    \draw[faintedge] (M-1-1) -- (M-3-2);

    \draw[faintedge] (M-1-2) -- (M-1-3);
    \draw[faintedge] (M-1-2) -- (M-2-2);
    \draw[faintedge] (M-1-2) -- (M-2-3);

    \draw[faintedge] (M-1-3) -- (M-2-3);

    \draw[faintedge] (M-2-1) -- (M-2-2);
    \draw[faintedge] (M-2-1) -- (M-3-1);
    \draw[faintedge] (M-2-1) -- (M-3-2);

    \draw[faintedge] (M-2-2) -- (M-2-3);
    \draw[faintedge] (M-2-2) -- (M-3-2);

    \draw[faintedge] (M-2-3) -- (M-3-3);
    \draw[faintedge] (M-3-1) -- (M-3-2);
    \draw[faintedge] (M-3-2) -- (M-3-3);

    \draw[faintedge] (M-1-1) to[bend left=20] (M-1-3);
    \draw[faintedge] (M-2-1) to[bend left=20] (M-2-3);
    \draw[faintedge] (M-3-1) to[bend left=20] (M-3-3);

    \draw[faintedge] (M-1-1) to[bend right=20] (M-3-1);
    \draw[faintedge] (M-1-2) to[bend right=20] (M-3-2);

    \draw[faintedge] (M-1-1) to[bend right=20] (M-3-3);
    \draw[faintedge] (M-1-2) -- (M-3-3);
    \draw[faintedge] (M-1-3) to[bend right=20] (M-3-3);

    \draw[faintedge] (M-2-1) -- (M-3-3);
    \draw[faintedge] (M-2-2) -- (M-3-3);

    % chosen subgraph edges (same adjacency)
    \draw[seledge] (M-2-1) -- (M-3-2);
    \draw[seledge] (M-2-1) -- (M-2-3);
    \draw[seledge] (M-1-2) -- (M-3-2); 
    \draw[seledge] (M-1-2) -- (M-2-3);

    % corner labels (optional, now reflecting shear)
    \node[lab, anchor=east] at ($(M-1-1)+(-0.25,0)$) {$(1.2,1.2)$};
    \node[lab, anchor=east] at ($(M-1-3)+(-0.25,0)$) {$(1.6,3.2)$};
    \node[lab, anchor=west] at ($(M-3-1)+(0.25,0)$) {$(3.2,1.6)$};
    \node[lab, anchor=west] at ($(M-3-3)+(0.25,0)$) {$(3.6,3.6)$};
\end{scope}
\end{tikzpicture}
\caption{The pertubation of $H(2143)$ in a $\CGrid_{3\times3}$.}
\label{fig:pertubed-cgrid}
\end{figure}